\section{Conclusiones}
Este trabajo se propuso desarrollar y evaluar un mecanismo automatizado que detectara la deriva arquitectónica combinando grafos de dependencia y \textit{Fitness Functions}, con el fin de clasificar el estado de conformidad estructural de proyectos de PyMES sin depender de documentación externa. Los resultados obtenidos permiten afirmar que el objetivo general se cumplió: el mecanismo fue construido, evaluado y demostró ser capaz de discriminar entre una implementación conforme y otra que acumula desviaciones, así como de detectar cada una de las tres familias de deriva cuando se presentan de forma aislada.

El cumplimiento de los objetivos específicos respalda esa afirmación. Respecto a OBJ1, se formuló un catálogo acotado de tres \textit{Fitness Functions} —integridad de capas, ausencia de ciclos y control de concentración de dependencias— en el que cada regla declara su propiedad protegida, su racionalidad, su unidad de verificación, su criterio de incumplimiento y su severidad; el catálogo no reside solo en el documento, sino que el propio prototipo lo deriva de la implementación vigente y puede imprimirlo bajo demanda. En cuanto a OBJ2, se implementó un componente de análisis estático que construye un grafo de dependencias entre paquetes a partir del \textit{bytecode}, prescindiendo por completo de la documentación arquitectónica y conservando, junto a cada arista, la referencia entre clases que la origina. En relación con OBJ3, se construyó un motor de validación que agrega las verificaciones en una Tasa de Conformidad y clasifica el estado del sistema en tres niveles interpretables, empleando los umbrales del $85\%$ y el $70\%$. Por último, OBJ4 se satisfizo mediante un experimento controlado sobre cinco escenarios en el que el mecanismo detectó el cien por ciento de las derivas inducidas sin producir un solo hallazgo sobre la implementación de referencia, superando el criterio de éxito del ochenta por ciento.

La tabla~\ref{tab:cierre} recoge esa correspondencia frente a los criterios de éxito que el anteproyecto fijó para cada objetivo, junto con la evidencia concreta que permite verificarla.

\begin{table}[H]
\centering
\caption{Cumplimiento de los criterios de éxito comprometidos en el anteproyecto.}
\label{tab:cierre}
\footnotesize
\begin{tabularx}{\textwidth}{p{1.1cm}p{5.0cm}X}
\toprule
\textbf{Obj.} & \textbf{Criterio de éxito comprometido} & \textbf{Evidencia de cumplimiento} \\
\midrule
OBJ1 & Catálogo con al menos una regla ejecutable por familia de deriva, con propósito, evidencia esperada y criterio de incumplimiento. &
Tres \textit{Fitness Functions} documentadas en la tabla~\ref{tab:catalogo}, derivadas de la implementación vigente e imprimibles desde el propio prototipo. \\
\midrule
OBJ2 & Grafo navegable que represente de forma consistente los componentes y dependencias necesarios para evaluar las reglas. &
Grafo construido desde el \textit{bytecode}, exportable en formato DOT con las capas agrupadas y las violaciones resaltadas; figuras~\ref{fig:grafo-base} y~\ref{fig:grafo-alterado}. \\
\midrule
OBJ3 & Cálculo automático de la CR, aplicación de los umbrales y clasificación coherente con la matriz de validación. &
Ecuación~\ref{eq:cr} y algoritmo~\ref{alg:conformidad}; clasificación coincidente con la matriz esperada en los cinco escenarios (tabla~\ref{tab:matrizObservada}). \\
\midrule
OBJ4 & Detección de al menos el $80\%$ de las derivas inducidas, línea base sin hallazgos inesperados y reportes comprensibles. &
Detección del $100\%$ de las seis derivas, cero hallazgos sobre la línea base y reportes con evidencia hasta la línea de código (tabla~\ref{tab:indicadores}). \\
\bottomrule
\end{tabularx}
\end{table}

El resultado más instructivo, sin embargo, no fue la confirmación de lo esperado sino el descubrimiento de una limitación. Al someter el mecanismo a escenarios con una única deriva —un diseño que el trabajo incorporó precisamente para medir la sensibilidad de cada regla por separado—, se constató que la Tasa de Conformidad, por ser una proporción no ponderada, diluye las violaciones localizadas hasta el punto de clasificar como aceptable un sistema que contiene una desviación real y documentada. La métrica mide con fidelidad la extensión del deterioro, pero es ciega a su gravedad, y esa ceguera se agrava conforme el sistema crece. La conclusión que de ello se desprende trasciende al prototipo: cualquier indicador de conformidad construido como proporción heredará esta propiedad, de modo que la agregación de hallazgos arquitectónicos no puede descansar únicamente en el conteo.

Los hallazgos dialogan con la literatura existente y la complementan. Frente a los enfoques basados en modelos, que exigen una especificación de referencia \citep{Murphy2001, SanMartin, KnodelPopescu2007}, la propuesta demuestra que es posible verificar la conformidad infiriendo la estructura directamente del código. Frente a los trabajos centrados en \textit{smells} \citep{Sas2019, SasAvgeriou}, la contribución consiste en traducir la detección en un veredicto agregado que sigue siendo reconstruible por quien lo recibe. Y frente a las propuestas de \textit{Fitness Functions} como práctica \citep{FordParsonsKua2017, PaulWang2019, ArchUnit}, este trabajo aporta un mecanismo que integra la extracción, la evaluación y la clasificación en un único instrumento ligero. En conjunto, la investigación confirma que la gobernanza arquitectónica continua no es un privilegio exclusivo de las grandes organizaciones, sino una capacidad que puede ponerse al alcance de las PyMES \citep{LAIG2025}.

\section{Contribuciones del trabajo}
Más allá del cumplimiento puntual de los objetivos, el trabajo aporta cuatro contribuciones que conviene deslindar.

La primera es \textbf{conceptual}: la propuesta articula en un único instrumento dos cuerpos de conocimiento que hasta ahora avanzaban por separado —el análisis estático de dependencias y las \textit{Fitness Functions} como práctica de gobernanza— y demuestra que su integración es viable y produce un diagnóstico accionable.

La segunda es \textbf{empírica} y no estaba prevista: la evidencia de que una tasa de conformidad formulada como proporción de verificaciones diluye las violaciones aisladas, junto con un criterio de clasificación que corrige esa ceguera combinando la extensión del incumplimiento con la gravedad del peor hallazgo. La corrección es deliberadamente simple —adoptar el más restrictivo de dos veredictos— y por ello sigue siendo explicable ante un equipo, que es la condición que el contexto de las PyMES impone a cualquier instrumento de gobernanza.

La tercera es \textbf{metodológica}: el documento ofrece un camino reproducible para construir un mecanismo de verificación de conformidad que no dependa de documentación externa, con decisiones de diseño explícitas, una matriz de validación fijada antes de la ejecución y una estrategia de verificación por etapas que separa el núcleo lógico de su integración con herramientas externas.

La cuarta es \textbf{práctica}: el prototipo constituye una base funcional que una PyME puede adoptar, adaptar y calibrar según su contexto, cuyo código está estructurado para favorecer la extensión y cuya evidencia se emite en formatos verificables por terceros.

Estas contribuciones dialogan con el estado del arte sin pretender superarlo en todos los frentes. La propuesta no compite con los enfoques basados en modelos en poder expresivo, ni con los enfoques de detección de \textit{smells} en profundidad analítica; su valor reside en el punto de equilibrio que logra entre costo, autonomía e interpretabilidad, precisamente el punto que los trabajos previos dejaban desatendido.

Esta lectura se alinea, sin desviarse, con la brecha que el anteproyecto identificó como espacio de contribución. Allí se constató que ningún antecedente reunía de forma simultánea cinco elementos: el análisis estático exclusivo del código, la combinación de grafo de dependencias con \textit{Fitness Functions}, la clasificación mediante una Tasa de Conformidad, un experimento controlado con criterios explícitos y una orientación específica a las PyMES sin documentación arquitectónica externa. El presente trabajo no añadió objetivos nuevos ni amplió el alcance: se limitó a ocupar, de manera verificable, ese espacio delimitado, y a documentar con honestidad dónde el instrumento resultó insuficiente y cómo se corrigió.

\section{Trabajos futuros}
La evaluación realizada abre varias líneas de trabajo que conviene explorar, ordenadas aquí por la urgencia que la propia evidencia les confiere.

La más apremiante es la \textbf{validación externa} sobre proyectos reales de la industria. Los escenarios utilizados son pequeños y sintéticos, y hay una razón concreta para esperar que la escala altere los resultados: la dilución de las violaciones localizadas se agrava conforme crece el número de verificaciones, de modo que en un sistema industrial el criterio basado únicamente en la Tasa de Conformidad sería aún menos discriminante que lo observado. Medir esa degradación con datos reales permitiría, además, validar el refinamiento sobre casos independientes de los que lo motivaron, que es la limitación de validez interna que este trabajo no logró contener.

Una segunda línea es la \textbf{calibración de los umbrales y el reconocimiento del rol de cada componente}. El autoanálisis del prototipo mostró que la regla de concentración, al operar solo sobre la topología, no distingue un componente dios de un concentrador legítimo como una raíz de composición o una fachada. Incorporar la posibilidad de declarar excepciones justificadas por configuración —y de registrarlas, para que la excepción sea también auditable— resolvería ese falso positivo conceptual sin renunciar a la regla.

Una tercera línea es el \textbf{análisis de la evolución temporal} de la conformidad. Si la métrica se registra en cada entrega, su serie histórica puede revelar tendencias de deterioro antes de que se conviertan en crisis, transformando el mecanismo de un detector puntual en un indicador prospectivo. Esta línea es especialmente prometedora porque la tendencia de una serie no padece la dilución que afecta al valor absoluto: un descenso sostenido es significativo aunque cada caída individual sea pequeña.

A ellas se suman tres extensiones de alcance más acotado: ampliar el catálogo de \textit{Fitness Functions} a otras manifestaciones de la deriva, como la violación de patrones de diseño; refinar la granularidad del análisis para emitir diagnósticos a nivel de clase y no solo de paquete; e integrar el mecanismo en una tubería de integración continua, para lo cual el prototipo ya ofrece la afordancia de terminar con un código de salida condicionado al veredicto. Por último, la extensión a otros ecosistemas mediante herramientas análogas a \textit{ArchUnit} en otras plataformas \citep{SugumarMedium} ampliaría el alcance de la propuesta más allá del ecosistema Java.

\section{Lecciones aprendidas}
El proceso dejó enseñanzas que trascienden el resultado puntual.

La primera, y la más valiosa, se refiere al diseño experimental: \textbf{un único escenario que acumula todas las derivas comprueba que el instrumento reacciona, pero oculta cómo reacciona}. Fue la decisión de construir escenarios con una sola deriva la que reveló la limitación de la métrica; con un solo caso alterado, el trabajo habría concluido que el mecanismo funcionaba correctamente y esa conclusión, aun siendo cierta en sus términos, habría sido engañosa. Aislar las variables no es un refinamiento estético del experimento: es lo que lo convierte en una fuente de conocimiento.

La segunda tiene que ver con la agregación de métricas. Reducir un diagnóstico a una sola cifra facilita la comunicación, pero toda agregación pierde información, y conviene saber cuál. En este caso, la proporción perdía la gravedad de los hallazgos, y recuperarla exigió añadir una segunda dimensión en lugar de complicar la fórmula de la primera. Mantener las dos lecturas separadas y visibles resultó preferible a fundirlas en un índice ponderado que nadie habría podido reconstruir.

La tercera se refiere a la configurabilidad: externalizar las reglas en un archivo de propiedades permitió reutilizar el mismo motor sobre seis sistemas distintos —los cinco escenarios y el propio prototipo— sin modificar una sola línea de código, una decisión de diseño que resultó decisiva.

La cuarta concierne al contrato de las reglas. Una primera versión permitía que los conteos que alimentan la métrica y la evidencia que la explica procedieran de recorridos distintos; obligar a que ambos surjan de una única evaluación eliminó por construcción una clase entera de inconsistencias posibles. En un instrumento cuya credibilidad descansa en que el número y su justificación coincidan, esa garantía no podía quedar en manos de la disciplina del programador.

La quinta se refiere a la validación por etapas —primero el núcleo lógico sobre grafos construidos a mano, después el componente que depende de la herramienta externa—, que redujo el riesgo de confundir un defecto propio con un comportamiento ajeno y permitió ejercitar casos límite que ningún caso de estudio ofrece. Aplicar además el instrumento sobre sí mismo resultó ser una comprobación barata y reveladora: confirmó que la arquitectura descrita era la implementada y, de paso, expuso una limitación de la regla de concentración que ningún escenario diseñado a propósito había mostrado.

Finalmente, una lección sobre el alcance. La disciplina de acotar el catálogo a tres dimensiones, en lugar de perseguir una cobertura exhaustiva, resultó esencial para entregar un instrumento funcional y evaluable. En investigación aplicada, la tentación de ampliar el alcance suele conspirar contra la posibilidad de validar; en este trabajo, la decisión inversa —acotar para poder verificar— fue la que dejó margen para investigar a fondo lo poco que se abarcó, y ese margen fue justamente donde apareció el hallazgo más interesante.
